\section{PEPS boundary tensors}

Let $d$ be the physical dimension. We allow horizontal virtual dimension
$D_{\mathrm h}$ and vertical virtual dimension $D_{\mathrm v}$ to differ. The
source uses a common virtual dimension, while the local contraction below also
makes sense with separate horizontal and vertical dimensions.

\begin{definition}[Local PEPS boundary tensor]
    \label{def:peps_local_boundary_tensor}
    \lean{TNLean.PEPS.localBoundaryMPOTensor}
    \leanok
    \uses{def:mpo_tensor}
    Let $A^s_{l,r,u,b}$ be a rank-five component function, with physical index
    $0\leq s<d$, horizontal virtual indices
    $0\leq l,r<D_{\mathrm h}$, and vertical virtual indices
    $0\leq u,b<D_{\mathrm v}$. The same component function is used at every
    lattice site. Contracting both its physical leg and its inward vertical leg
    between a ket copy and a bra copy defines an MPO tensor $E_A$ with physical
    dimension $D_{\mathrm v}$ and bond dimension $D_{\mathrm h}^2$.

    The unlabelled diagram at source lines~721--724 of
    \cite{Cirac2016MPDO_arXiv} depicts this local ket-bra boundary tensor. The
    choice of the lower leg as the inward ket leg fixes the orientation of the
    displayed component formula. We label the bra legs after reflecting the bra
    copy into the ket coordinate frame. Thus the contracted leg is called $b$
    in both copies; before reflection, it is the ket-down leg joined to the
    bra-up leg in the stacked diagram.
\end{definition}

\begin{theorem}[Component formula for the local boundary tensor]
    \label{thm:peps_local_boundary_tensor_formula}
    \lean{TNLean.PEPS.localBoundaryMPOTensor_apply,
        TNLean.PEPS.localBoundaryMPOTensor_finProdFinEquiv_apply}
    \leanok
    \uses{def:peps_local_boundary_tensor}
    With $b$ denoting the contracted inward vertical leg,
    \begin{align}
        (E_A)^{u,u'}_{(l,l'),(r,r')}
         & =\sum_{s=0}^{d-1}\sum_{b=0}^{D_{\mathrm v}-1}
        A^s_{l,r,u,b}
        \overline{A^s_{l',r',u',b}}.
        \label{eq:peps_local_boundary_tensor_component}
    \end{align}
    The physical coordinates are the remaining ket and bra indices $u,u'$, and
    the bond coordinates are $(l,l')$ and $(r,r')$. We order the five legs as
    left, right, up, down, physical and pack each horizontal ket-bra pair into
    one product index, using the reflected bra coordinate convention from the
    definition.
\end{theorem}

\begin{proof}
    \leanok
    \uses{def:peps_local_boundary_tensor}
    Expand the contractions over the common physical and inward vertical
    indices, and decode the two packed bond coordinates. This gives
    (\ref{eq:peps_local_boundary_tensor_component}).
\end{proof}

\begin{definition}[Boundary fixed point]
    \label{def:peps_transfer_fixed_point}
    \notready
    \discussion{7371}
    The row-to-row transfer operator of a homogeneous PEPS is represented by a
    strip of local ket-bra doubled tensors. A boundary fixed point is a left or
    right boundary operator invariant under the corresponding transfer action.
    The transfer-operator diagram at source lines~709--712 of
    \cite{Cirac2016MPDO_arXiv} depicts the operator from which the fixed point is
    selected; it does not depict the fixed point itself. An algebraic definition
    must specify the transfer space, orientation, and the choice of left and
    right fixed points.
\end{definition}

\begin{definition}[Algebraic PEPS renormalization fixed-point relation]
    \label{def:peps_rfp_relation}
    \notready
    \discussion{7371}
    The PEPS renormalization diagram at source lines~716--720 of
    \cite{Cirac2016MPDO_arXiv} identifies two neighboring tensors with one
    effective local tensor, up to an isometry on their physical indices. An
    algebraic definition must specify the domain and codomain of that isometry
    and the identifications of the external virtual spaces.
\end{definition}

\begin{definition}[Boundary coarse-graining channel]
    \label{def:peps_boundary_coarse_graining}
    \notready
    \discussion{7371}
    \uses{def:peps_transfer_fixed_point, def:peps_rfp_relation,
        def:peps_local_boundary_tensor, def:is_rfp_via_ts}
    A coarse-graining channel $\mathcal S$ acts from two boundary sites to one
    and satisfies $\mathcal S[M_2(X)]=M_1(X)$ for every virtual operator $X$.
    The coarse-graining map in the diagram at source lines~725--729 of
    \cite{Cirac2016MPDO_arXiv} must be defined on the full two-site boundary
    algebra and proved completely positive and trace-preserving, not merely on
    the image selected by the PEPS isometry.
\end{definition}

\begin{definition}[Boundary refinement channel]
    \label{def:peps_boundary_refinement}
    \notready
    \discussion{7371}
    \uses{def:peps_transfer_fixed_point, def:peps_rfp_relation,
        def:peps_local_boundary_tensor, def:is_rfp_via_ts}
    A refinement channel $\mathcal T$ acts from one boundary site to two and
    satisfies $\mathcal T[M_1(X)]=M_2(X)$ for every virtual operator $X$. It
    is obtained in the diagram at source lines~725--729 of
    \cite{Cirac2016MPDO_arXiv} by inserting the displayed local factor. The
    factor and its normalization must be specified so that the map on the full
    one-site boundary algebra is completely positive and trace-preserving.
\end{definition}

\begin{theorem}[PEPS renormalization induces a boundary MPO fixed point]
    \label{thm:peps_rfp_implies_boundary_rfp}
    \notready
    \discussion{7371}
    \uses{def:peps_transfer_fixed_point, def:peps_rfp_relation,
        def:peps_boundary_coarse_graining, def:peps_boundary_refinement,
        def:is_rfp_via_ts}
    Suppose a homogeneous PEPS satisfies the neighboring blocking relation at
    source lines~716--720 of \cite{Cirac2016MPDO_arXiv}, and let its transfer
    fixed points determine the boundary MPO displayed at source lines~721--724.
    Then that boundary MPO is a renormalization fixed point in the sense of
    Definition~\ref{def:is_rfp_via_ts}; the maps $\mathcal S$ and $\mathcal T$
    are those displayed at source lines~725--729.
\end{theorem}

\begin{proof}
    \notready
    \uses{def:peps_boundary_coarse_graining,
        def:peps_boundary_refinement, def:is_rfp_via_ts}
    The source diagrams do not specify the domain and codomain of the physical
    isometry, the virtual-space identifications, the choice and normalization of
    the transfer fixed points, or the inserted factor in the refinement map.
    Consequently, complete positivity, trace preservation, and the identities
    $\mathcal S[M_2(X)]=M_1(X)$ and
    $\mathcal T[M_1(X)]=M_2(X)$ cannot yet be checked algebraically. This
    missing specification is recorded in
    \cite{gap:cpsv16_peps_boundary_rfp_specification}.
\end{proof}
